	% --- CONFIGURAÇÕES DE PACOTES ---
%Comando pra conseguir identificar facilmente textos para correção
\newcommand{\pendente}[1]{\textcolor{red}{#1}}
% Configurações do pacote backref
\renewcommand{\backrefpagesname}{Citado na(s) página(s):~}
\renewcommand{\backref}{}
\renewcommand*{\backrefalt}[4]{
	\ifcase #1 %
	Nenhuma citação no texto.%
	\or
	Citado na página #2.%
	\else
	Citado #1 vezes nas páginas #2.%
	\fi}

% Configurações de aparência do PDF final
\definecolor{blue}{RGB}{41, 5, 195}

\makeatletter
\hypersetup{
	pdftitle={\@title}, 
	pdfauthor={\@author},
	pdfsubject={\imprimirpreambulo},
	pdfcreator={LaTeX with abnTeX2},
	pdfkeywords={abnt}{latex}{abntex}{abntex2}{trabalho acadêmico}, 
	colorlinks=true,       		
	linkcolor=black,     		
	citecolor=black,     		
	filecolor=black,     		
	urlcolor=black,      		
	bookmarksdepth=4
}
\makeatother

% Espaçamentos entre linhas e parágrafos 
\setlength{\parindent}{1.3cm}
\setlength{\parskip}{0.2cm} 

% Ajustes de hifenização e quebra de linha
\sloppy
\clubpenalty=10000
\widowpenalty=10000
\hyphenpenalty=10000

% Compila o indice
\makeindex

%Comando pra dar espaço na parte da reação de transesterificação
\newcommand{\etapaspace}{\vspace{1.5em}}